\clearpage\begingroup\SourceLeaf{107b}\HistoricalHead{102}{L}{PHYSIK. BRIEFWECHSEL.}
\begin{center}[19.]\end{center}
\begin{center}\textsc{Gauss} an \textsc{Encke}.\end{center}
\noindent .\ .\ .\ .\ .\ .\ .\ .\ .\ .\ .\ .\ .\ .\ .\ .\ .\ .\ .\ .\ .\par
\indent Ich bin nicht abgeneigt, künftig eine Schrift etwa \textit{Zum Gebrauch des
Magnetometers} abzufassen, welche ohne sich in tiefe Theorie einzulassen, alle
für den Liebhaber und den feinern Beobachter erforderliche Anleitung enthielte.
Vielleicht wäre eine Briefform am angemessensten. Läugnen kann
ich freilich nicht, dass dergleichen populäre Schriftstellerei nicht gerade meine
Liebhaberei ist, dass ich sie vielmehr als ein Opfer ansehe, welches zu bringen
nur die Nützlichkeit mich bewegen kann, zumal da mir soetwas, nach dem
ersten Versuch dieser Art, welchen ich in \textsc{Schumachers} Jahrbuch gemacht
\clearpage\SourceLeaf{108}\HistoricalHead{103}{R}{GAUSS AN ENCKE.}
\noindent habe\,[\textsuperscript{*})], viel mehr Zeit kostet, als Sie ihm vielleicht ansehen. Ich werde
nun zuerst erwarten, ob unser Publicum solche Opfer zu erkennen weiss.\begingroup\renewcommand{\thefootnote}{}
\footnotetext{[\textsuperscript{*}) \textit{Erdmagnetismus und Magnetometer}, Jahrbuch für 1836, herausgegeben von \textsc{Schumacher}, Stuttgart u. Tübingen 1836, S. 1; Werke V, S. 315.]}
\endgroup

Mehr sagen meinem Geschmack \textls[120]{zwei} andere damit nahe zusammenhängende
Arbeiten zu, wovon auch bereits einiges niedergeschrieben ist. Das
eine, über allgemeine Theoreme die im verkehrten Verhältniss des Quadrats
der Entfernungen wirkende Kräfte betreffend, wo u. a. das Theorem (\textit{Intensitas
etc.} p. [7\textsuperscript{**})]) bewiesen werden würde; das zweite über die Grundgesetze der
Galvanischen Ströme und der Induction und deren Zurückführung auf absolute
Maasse.\begingroup\renewcommand{\thefootnote}{}
\footnotetext{[\textsuperscript{**}) Siehe \textit{Intensitas vis magneticae terrestris ad mensuram absolutam revocata}, art. 2., Werke V,
S. 87. Es handelt sich um den folgenden Satz: »Wie immer auch die Verteilung des freien Magnetismus
im Innern eines Körpers geartet sei, so kann sie nach einem bestimmten Gesetze durch eine andere Verteilung
allein auf der Oberfläche des Körpers ersetzt werden, die in bezug auf die nach Außen hin wirkenden
Kräfte jener genau gleichwertig ist, so daß jedes außerhalb des Körpers befindliche Element magnetischer
Flüssigkeit von der im Innern des Körpers wirklich vorhandenen Verteilung dieselben anziehenden
oder abstoßenden Kräfte erfährt, wie von der gedachten Verteilung auf der Oberfläche«. Ein Beweis dieses
Satzes findet sich in der Abhandlung vom Jahre 1839 \textit{Allgemeine Lehrsätze in Beziehung auf die im verkehrten
Verhältnisse des Quadrats der Entfernung wirkenden Anziehungs- und Abstossungskräfte}, Art. 36,
Werke V, S. 241.]}
\endgroup\ Beides sind aber keine Arbeiten, die man nur aus dem Aermel schütteln
könnte. Ich könnte noch eine vierte Arbeit nennen, nemlich die Allgemeine
Theorie des Erdmagnetismus, die ich auszuführen mich anheischig machen
wollte, wenn die Thatsachen reichlich genug und in der \textls[120]{nothwendigen
Form} vorlägen. Diese Form muss aber eine ganz andere sein, als die gewöhnliche,
wonach die drei Elemente so gefasst werden: Declination, Inclination
und \textls[120]{ganze Intensität}. Diese Form ist für meinen Zweck ganz unbrauchbar.
Ich würde aber die Arbeit unternehmen, wenn wir nur noch eine Intensitäts-Karte
hätten, die die \textls[120]{horizontale Intensität} allein darstellte, obwohl auch
dann erst noch eine Umschmelzung nöthig wäre. \textsc{Hansteen} hatte mir Hoffnung
gemacht, mir die drei Karten in der gewünschten Form zu liefern,
scheint aber sein Versprechen vergessen zu haben.

Schon vor einigen Jahren hatte ich eine besondere Methode erdacht, die
Inclination zu bestimmen, wozu eine Nadel mit einer feinen Axe etwas ausserhalb
des Schwerpunktes nöthig ist, und wo man bloss Schwingungsdauern beobachtet.
Ein damals von mir gemachter ganz roher Versuch liess die Brauch-
\clearpage\SourceLeaf{109}\HistoricalHead{104}{L}{PHYSIK. BRIEFWECHSEL.}
\noindent barkeit dieser Methode schon erkennen, aber die Herren \textsc{Sartorius} und \textsc{Listing}
haben sie mit einem besser gearbeiteten Apparat in Neapel recht glücklich
ausgeführt\,[\textsuperscript{*})], sodass ich diese Methode für die allergenaueste halten möchte.\begingroup\renewcommand{\thefootnote}{}
\footnotetext{[\textsuperscript{*}) Vergl. Resultate aus den Beobachtungen des magnetischen Vereins im Jahre 1838, S. 58.]}
\endgroup\ Ihr Resultat ist aber noch nicht definitiv, da sie ein Element der Nadel noch
nicht scharf bestimmt haben, was sie sich vorbehalten, nach ihrer Rückkehr
zu thun. Auch absolute Declination und Intensität haben sie mit vieler Sorgfalt
bestimmt; der erste Versuch nach Göttingen. Den dritten wird vermuthlich
\textsc{Svanberg} machen, falls Sie ihm nicht zuvorkommen.

Auch hier in Göttingen ist im Laufe des verflossenen Jahres manches
Neue experimentirt. Wir haben die Reibungselectricität durch die ganze
Kette vom physikalischen Cabinet zur Sternwarte getrieben, ihre magnetische
Wirkung gemessen, und gefunden, dass unter gewissen Vorsichtsregeln nur
wenig unterwegs verloren geht. Die magnetische Telegraphie durch Inductionsstösse
haben wir oft so ausgeführt, dass 7 und mehr Buchstaben in der Minute
transmittirt werden. Dass anstatt des früheren Inductors von 1050 Umwindungen
einer von 3537 angeschafft ist, haben Sie schon aus dem oben angeführten
Aufsatz [Werke V, S. 340] ersehen. Im Sept[ember] erhielten wir zuerst damit
physiologische Wirkungen, von denen Ihnen vielleicht H[err] \textsc{Poggendorf} erzählt
hat, da er gerade bei den allerersten Versuchen gegenwärtig war und selbst
Theil daran nahm. Aber was ich erst später fand, hat wie mir deucht eine
grosse Merkwürdigkeit, nemlich dass man, wenn man den Strom mit den Lippen
auffängt, d. i. an einer Lippe den Strom einströmen, an der anderen ausströmen
lässt, man die Richtung des Stromes jedesmahl bestimmt unterscheidet.
Immer empfindet man nemlich an der Lippe, wo der negative einströmt,
\textls[120]{allein}, wenn der Strom nicht zu stark ist, oder im entgegengesetzten Fall
doch ganz überwiegend. Dieser Versuch, den wir viele hundertmahl gemacht,
und selbst scherzweise zum Telegraphieren gebraucht haben, schlägt nie fehl.

Seit ein Paar Monaten habe ich nun aber den Inductor abermals fast
auf das doppelte verstärken lassen, sodass er gegen 7000 Umwindungen hat.
Hier sind die physiologischen Wirkungen weit stärker; man kann den Strom
selbst durch die nur wenig befeuchteten \textls[120]{Hände} gehend sehr fühlbar machen,
was bei dem schwächern Inductor nicht anging. Das allermerkwürdigste
\clearpage\SourceLeaf{110}\HistoricalHead{105}{R}{GAUSS AN ENCKE.}\HistoricalFoot{XI1.}{14}
\noindent aber scheint zu sein, dass (seit etwa 8—10 Tagen) es mir gelingt Funken
hervorzubringen, und dass diese Funken, nach der Richtung des Stromes eine
verschiedene Farbe haben. Ich lasse die Funken von einer feinen Nadel
gegen eine \textls[120]{feste} Metallfläche (nicht Quecksilber) überschlagen, wo der positive
Strom einen hellgelben Funken mit einem Stich ins \textls[120]{Grüne} giebt, der negative
einen \textls[120]{violetten} oder zuweilen \textls[120]{roth}gelben Funken. Mehreremale hat der letztere
die Spitze der Nadel geschmolzen. Die Entfernung muss dabei aber sehr
klein sein, vermuthlich unter \(\frac{1}{1000}\) Zoll, und ich bediene mich zur Stellung
einer sehr feinen Schraube (etwa 80 Gewinde auf den Zoll). Mit einer anderen
viel grobern Schraube hat der Versuch bisher noch nicht gelingen wollen.
Die Funken erfolgen selbst, wenn der Strom durch die \textls[120]{ganze} Kette geht,
jetzt fast eine Meile lang. Es werden sich hieran noch eine Menge höchst
interessanter Versuche knüpfen lassen, besonders wenn ich meine Vorrichtungen
erst etwas bequemer eingerichtet habe.

Herzlich wünsche ich Ihnen Glück zur Vollendung Ihrer Sternwarte. Der
Besitz so reicher Hilfsmittel muss Ihnen grossen Genuss gewähren.

Stets mit unwandelbarer Freundschaft
\begin{flushright}der Ihrige\par\medskip C. F. \textsc{Gauss}.\end{flushright}
\par\indent Göttingen den [\textsuperscript{*})] Januar 1836.\begingroup\renewcommand{\thefootnote}{}
\footnotetext{[\textsuperscript{*}) Der Tag ist in der Handschrift nicht angegeben.]}
\endgroup\par
\indent Herrn von \textsc{Humboldt}, der jetzt vermuthlich nach Berlin zurückgekehrt
sein wird, bitte ich angelegentlich mich bestens zu empfehlen. In den letzten
Tagen haben wir auch angefangen, uns mit den Thermogalvanischen Strömen
zu beschäftigen. Vermittelst einer einfachen Vorrichtung können wir solche
durch die ganze 1 Meile lange Kette treiben, sodass z. B. von der Sternwarte
aus sämmtliche 4 Magnetometer (worunter jetzt zwei mit 25 pf.Stäben, da
\textsc{Weber} sich auch einen solchen aufgehängt hat) in sehr bedeutende Bewegungen
gesetzt werden können. Die Natur dieser Ströme ist durchaus nicht
von der auf andere Art hervorgebrachten verschieden.
\clearpage\endgroup
